% !PRINTED 16
\weberpage{16}
%<<CONT
Zahl $\alpha$ durch sehr verschiedene Zahlenreihen erzeugt werden.
Diese Zahlenreihen sind aber alle als unter einander gleich zu
betrachten. Man kann unter Anderem die Zahlen von $S$ alle als
rationale Brüche annehmen. Es lässt sich auch umgekehrt leicht
nachweisen, dass man zu jeder gegebenen Zahl $\alpha$ immer
Zahlenreihen $S$ angeben kann, durch die $\alpha$ erzeugt wird, so dass
also die Gesammtheit der Zahlenreihen $S$ gleichfalls eine stetige
Menge bildet.

\begin{center}
\rule{0.18\linewidth}{0.4pt}
\end{center}

Bei der Erklärung der Grundrechnungsarten hat sich bei der Subtraction
eine unbequeme Beschränkung ergeben, von der wir uns freimachen durch
Einführung der \gesperrt{Null und der negativen Zahlen}.

Es möge $x$ jedes Element des bisher definirten Zahlensystems
bedeuten, das wir jetzt als das System der \gesperrt{positiven Zahlen}
bezeichnen wollen. Wir nehmen dies Zahlensystem ein zweites Mal und
bezeichnen zum Unterschied in diesem zweiten System, das als das
System der \gesperrt{negativen Zahlen} bezeichnet werden soll, jedes
Element mit $-x$. Das zweite System ordnen wir nun dem ersten gerade
entgegengesetzt, so dass überall, wo in dem System $x$ „grösser“
steht, in dem System $-x$ „kleiner“ gesetzt wird, und umgekehrt.
Addition und Subtraction werden in $-x$ ebenso erklärt wie in $x$, so
dass $(-x)+(-y)=-(x+y)$, $(-x)-(-y)=-(x-y)$ sein soll.

Wir wollen aber diese beiden Zahlensysteme in der Weise zusammenordnen,
dass jedes $-x$ kleiner sein soll als jedes $x$. Wir erhalten so eine
geordnete Menge, in der kein grösstes und kein kleinstes Element
vorhanden ist. Dieses System ist auch im Allgemeinen stetig, nur der
einzige Schnitt $(-x,x)$ wird durch kein Element erzeugt, und hier, wo
beide Systeme zusammenstossen, ist also noch eine Verletzung der
Stetigkeit vorhanden. Um die Stetigkeit herzustellen, müssen wir also
dem Schnitt $(-x,x)$ entsprechend noch eine Zahl \gesperrt{Null oder 0}
hinzufügen, die eben durch diesen Schnitt definirt ist. Dann haben wir
eine geordnete stetige, beiderseits unbegrenzte Menge, die vollständige
Reihe der \gesperrt{reellen Zahlen}.

In dem so erweiterten Zahlenbereiche erklären wir nun die Addition
allgemein, indem wir definitionsweise setzen:
%CONT>>
